%% kapa: exam=1 round=35 year=2024 period=1 form=A subject=economics subject_ko=경제학원론 seq=12 number=52
%% source: Q-Net 감정평가사 기출문제 (cst00309) artlSeq=5213229 file=raw/1008/5213229/1교시 문제지 원본.pdf
%% conversion: kordoc 4.24.51 -> tools/gen1.py (auto); review=proofread-sonnet; flags=none
\begin{question}{52}
정부의 실효성 있는 가격규제의 효과에 관한 설명으로 옳은 것은? (단, 수요곡선은 우하향, 공급곡선은 우상향한다.)
\begin{choices}
  \item 가격상한제가 실시되면, 시장에서의 실제 거래량은 실시 이전보다 증가할 것이다.
  \item 가격하한제가 실시되면, 시장에서의 실제 거래량은 실시 이전보다 증가할 것이다.
  \item 최저임금제는 가격상한제에 해당하는 가격규제이다.
  \item 가격하한제가 실시되면, 초과수요가 발생하여 암시장이 형성된다.
  \item 가격상한제와 가격하한제 모두 자중손실(deadweight loss)이 발생한다.
\end{choices}
\end{question}
\answer{5}
